\documentclass[12pt]{article}
\usepackage[utf8]{inputenc}
\usepackage[T1]{fontenc}
\usepackage{hyperref}
\title{Izere Language Resources for K-12 Curriculum Development}
\author{Compiled by Claude Sonnet 5.5 (claude-sonnet-5-5) for Andr\'as Kornai}
\date{}
\begin{document}
\maketitle

\section{Introduction}
The code fiz is a retired ISO 639-3 code formerly assigned to Izere, a Plateau dialect continuum of Nigeria (mostly Kaduna, Bauchi and Plateau states). It is now covered by active codes such as izr (Izere), cen (Cèn) and gne (Ganàng). The ``family: Bookkeeping'' entry in the task identification looks like an artifact of the code's retired status and is not a real language family. Glottolog treats Izere as one language (izr), classed as threatened. This report is based on a single search and is very short.

\section{Primary Educational Resources}
None found.

\section{Linguistic and Reference Materials}
\begin{itemize}
\item Wikipedia, Izere language: \url{https://en.wikipedia.org/wiki/Izere_language}
\item Glottolog entry for izr: \url{https://glottolog.org/resource/languoid/iso/izr}
\item PHOIBLE inventory entry for izr: \url{https://phoible.org/languages/izr}
\end{itemize}

\section{Religious and Community Resources}
None found.

\section{Digital and Audio Resources}
None found.

\section{Curriculum Development Context}
Work on this language should proceed under the active code izr (and cen, gne) rather than fiz. Not to be confused with Idere (ide), an unrelated language of Akwa Ibom State. No school materials were found in this search; further searching under izr is recommended.

\end{document}